\begin{thebibliography}{99}
\addcontentsline{toc}{section}{Список литературы}

\bibitem{sohl2024}
\emph{Sohl-Dickstein J.} The boundary of neural network trainability is fractal
// arXiv preprint arXiv:2402.06184. --- 2024. --- URL:
\url{https://arxiv.org/abs/2402.06184} (дата обращения: 20.05.2026).

\bibitem{torkamandi2025}
\emph{Torkamandi B. [и др.]} Mapping the edge of chaos: fractal-like boundaries
in the trainability of decoder-only transformer models // arXiv preprint
arXiv:2501.04286. --- 2025. --- URL: \url{https://arxiv.org/abs/2501.04286}
(дата обращения: 20.05.2026).

\bibitem{liu2024}
\emph{Liu Y., Arnould L., Sohl-Dickstein J.} Complex fractal trainability
boundary can arise from trivial non-convexity // arXiv preprint
arXiv:2406.13971. --- 2024. --- URL: \url{https://arxiv.org/abs/2406.13971}
(дата обращения: 20.05.2026).

\bibitem{cohen2021}
\emph{Cohen J. M., Kaur S., Li Y., Kolter J. Z., Talwalkar A.} Gradient descent
on neural networks typically occurs at the edge of stability //
International Conference on Learning Representations (ICLR). --- 2021. --- URL:
\url{https://arxiv.org/abs/2103.00065} (дата обращения: 20.05.2026).

\bibitem{cohen2022}
\emph{Cohen J. M. [и др.]} Adaptive gradient methods at the edge of stability
// arXiv preprint arXiv:2207.14484. --- 2022. --- URL:
\url{https://arxiv.org/abs/2207.14484} (дата обращения: 20.05.2026).

\bibitem{lewkowycz2020}
\emph{Lewkowycz A., Bahri Y., Dyer E., Sohl-Dickstein J., Gur-Ari G.} The large
learning rate phase of deep learning: the catapult mechanism // arXiv preprint
arXiv:2003.02218. --- 2020. --- URL: \url{https://arxiv.org/abs/2003.02218}
(дата обращения: 20.05.2026).

\bibitem{kong2020}
\emph{Kong L., Tao M.} Stochasticity of deterministic gradient descent: large
learning rate for multiscale objective function // Advances in Neural
Information Processing Systems (NeurIPS). --- 2020. --- URL:
\url{https://proceedings.neurips.cc/paper/2020/file/1b9a80606d74d3da6db2f1274557e644-Paper.pdf}
(дата обращения: 20.05.2026).

\bibitem{chen2023}
\emph{Chen X., Balasubramanian K., Ghosal P., Agrawalla B.} From stability to
chaos: analyzing gradient descent dynamics in quadratic regression //
arXiv preprint arXiv:2310.01687. --- 2023. --- URL:
\url{https://arxiv.org/abs/2310.01687} (дата обращения: 20.05.2026).

\bibitem{feigenbaum1978}
\emph{Feigenbaum M. J.} Quantitative universality for a class of nonlinear
transformations // Journal of Statistical Physics. --- 1978. --- Т. 19,
№ 1. --- С. 25--52. --- DOI: 10.1007/BF01020332.

\bibitem{liyorke1975}
\emph{Li T.-Y., Yorke J. A.} Period three implies chaos // The American
Mathematical Monthly. --- 1975. --- Т. 82, № 10. --- С. 985--992. ---
DOI: 10.2307/2318254.

\bibitem{polyak1964}
\emph{Поляк Б. Т.} О некоторых способах ускорения сходимости итерационных
методов // Журнал вычислительной математики и математической физики. ---
1964. --- Т. 4, № 5. --- С. 791--803.

\bibitem{nesterov1983}
\emph{Нестеров Ю. Е.} Метод решения задачи выпуклого программирования со
скоростью сходимости $O(1/k^2)$ // Доклады АН СССР. --- 1983. --- Т. 269,
№ 3. --- С. 543--547.

\bibitem{kingma2015}
\emph{Kingma D. P., Ba J.} Adam: a method for stochastic optimization //
International Conference on Learning Representations (ICLR). --- 2015. --- URL:
\url{https://arxiv.org/abs/1412.6980} (дата обращения: 20.05.2026).

\bibitem{tieleman2012}
\emph{Tieleman T., Hinton G.} Lecture 6.5---RMSprop: divide the gradient by a
running average of its recent magnitude // COURSERA: Neural Networks for
Machine Learning. --- 2012.

\bibitem{loshchilov2019}
\emph{Loshchilov I., Hutter F.} Decoupled weight decay regularization //
International Conference on Learning Representations (ICLR). --- 2019. --- URL:
\url{https://arxiv.org/abs/1711.05101} (дата обращения: 20.05.2026).

\bibitem{bottou2018}
\emph{Bottou L., Curtis F. E., Nocedal J.} Optimization methods for large-scale
machine learning // SIAM Review. --- 2018. --- Т. 60, № 2. --- С. 223--311. ---
DOI: 10.1137/16M1080173.

\bibitem{smith2018}
\emph{Smith S. L., Kindermans P.-J., Ying C., Le Q. V.} Don't decay the learning
rate, increase the batch size // International Conference on Learning
Representations (ICLR). --- 2018. --- URL:
\url{https://arxiv.org/abs/1711.00489} (дата обращения: 20.05.2026).

\bibitem{goodfellow2016}
\emph{Goodfellow I., Bengio Y., Courville A.} Deep Learning. --- Cambridge, MA:
MIT Press, 2016. --- 800 с. --- URL: \url{http://www.deeplearningbook.org}
(дата обращения: 20.05.2026).

\bibitem{falconer2014}
\emph{Falconer K.} Fractal Geometry: Mathematical Foundations and
Applications. --- 3rd ed. --- Chichester: Wiley, 2014. --- 400 с.

\bibitem{mandelbrot1982}
\emph{Mandelbrot B. B.} The Fractal Geometry of Nature. --- New York:
W. H. Freeman and Company, 1982. --- 468 с.

\bibitem{strogatz2015}
\emph{Strogatz S. H.} Nonlinear Dynamics and Chaos: With Applications to
Physics, Biology, Chemistry, and Engineering. --- 2nd ed. --- Boulder, CO:
Westview Press, 2015. --- 513 с.

\bibitem{porespy2019}
\emph{Gostick J. T. [и др.]} PoreSpy: a Python toolkit for quantitative
analysis of porous media images // Journal of Open Source Software. ---
2019. --- Т. 4, № 37. --- С. 1296. --- URL:
\url{https://porespy.org/examples/metrics/tutorials/computing_fractal_dim.html}
(дата обращения: 20.05.2026).

\bibitem{harris2020}
\emph{Harris C. R. [и др.]} Array programming with NumPy // Nature. ---
2020. --- Т. 585. --- С. 357--362. --- DOI: 10.1038/s41586-020-2649-2.

\bibitem{virtanen2020}
\emph{Virtanen P. [и др.]} SciPy 1.0: fundamental algorithms for scientific
computing in Python // Nature Methods. --- 2020. --- Т. 17. --- С. 261--272. ---
DOI: 10.1038/s41592-019-0686-2.

\bibitem{pedregosa2011}
\emph{Pedregosa F. [и др.]} Scikit-learn: machine learning in Python //
Journal of Machine Learning Research. --- 2011. --- Т. 12. --- С. 2825--2830.

\bibitem{skimage2014}
\emph{van der Walt S. [и др.]} scikit-image: image processing in Python //
PeerJ. --- 2014. --- Т. 2. --- С. e453. --- DOI: 10.7717/peerj.453.

\end{thebibliography}
